\documentclass[10pt,a4paper]{amsart}
\usepackage[margin=19mm]{geometry}
\usepackage{amsmath,amssymb,mathtools}
\usepackage[scr=rsfs,cal=euler]{mathalfa}
\usepackage{xfrac}
\usepackage[unicode,hidelinks,bookmarks=false]{hyperref}
\newcommand{\ie}{i.e.\ }
\newcommand{\cf}{cf.\ }
\newcommand{\resp}{respectively\ }
\newcommand{\Subsection}{\subsection}
\providecommand{\nobreakdash}{\nobreak\hbox{-}}
\setlength{\emergencystretch}{2em}
\multlinegap0pt
\usepackage{bm}
\usepackage{bbm}
\usepackage{dsfont}
\usepackage{xspace}
\usepackage{cancel}
\usepackage{mathtools}
\usepackage{amsthm}
\makeatletter
\@ifundefined{theorem}{\newtheorem{theorem}{Theorem}}{}
\@ifundefined{proposition}{\newtheorem{proposition}[theorem]{Proposition}}{}
\makeatother
\DeclareMathOperator{\coker}{coker}
\newcommand{\C}{\mathbb{C}}
\newcommand{\M}{\mathcal{M}}
\newcommand{\N}{\mathcal{N}}
\newcommand{\Z}{\mathbb{Z}}
\newcommand{\ch}{\mathrm{ch}}
\newcommand{\g}{\mathfrak{g}}
\newcommand{\h}{\mathfrak{h}}
\newcommand{\xto}[1]{\xrightarrow{#1}}
\renewcommand{\P}{\mathcal{P}}
\renewcommand{\sl}{\mathfrak{sl}}
\title[A formula for the $q$-character of functions on the nilpoten]{A formula for the $q$-character of functions on the nilpotent cone of some Lie algebra representations}
\author{Vasily Krylov, Frank Wang}
\date{Source: arXiv:2608.03314v1 (CC BY 4.0)}
\begin{document}
\maketitle

\noindent\emph{Source and licence.} A formula for the $q$-character of functions on the nilpotent cone of some Lie algebra representations. Version arXiv:2608.03314v1 (CC BY 4.0), distributed under the Creative Commons Attribution 4.0 International licence, as recorded in the arXiv licence metadata for this exact version and shown on the abstract page https://arxiv.org/abs/2608.03314v1. Full rights notice: licenses/arxiv-2608.03314-CC-BY-4.0.txt.\par\smallskip

\noindent\textit{Editorial scope.} Counted unit: the Proposition whose source label is \texttt{prop:2cycle} in the article (source0.tex lines 551--553, source label \texttt{prop:2cycle}), delivered together with the article's own complete original proof of it (source0.tex lines 555--599, one \texttt{proof} environment of 45 lines). No mathematics is written, repaired or re-derived: statement and proof are the article's own text line for line. Required context retained uncounted: thm (lines 126--129); prop:hesselink (lines 190--199). Every definition, display and label of those lines is retained unchanged. Omitted from this package and not redistributed: 1--125, 130--189, 200--550, 554--554, 600--679. Nothing inside a retained excerpt is omitted or altered: every retained body is delivered exactly as the article's lines are written, so no omission receipt of any kind accompanies this package. Citations: the retained excerpts carry no citation command at all, so no bibliography entry of the article is transcribed and no reference is redistributed. Reference adaptation: none was needed and none was made; every reference inside the delivered unit resolves inside it, so no unresolved reference or citation is produced. Printed slips kept literal: no printed word, symbol, hypothesis or number of the counted statement or of its proof was altered. Layout and class: the article's own class is \texttt{article} with packages including inputenc, amsmath, amssymb, amsthm, enumitem, geometry, setspace, graphicx, hyperref, tikz, tikz-cd, float, verbatim, biblatex; the wrapper is the lane's pinned \texttt{amsart} editorial wrapper at 10pt on A4, which restates the theorem-like environments the retained text needs and adds the article's own macro definitions that the retained text needs (\texttt{\textbackslash C}, \texttt{\textbackslash M}, \texttt{\textbackslash N}, \texttt{\textbackslash P}, \texttt{\textbackslash Z}, \texttt{\textbackslash ch}, \texttt{\textbackslash coker}, \texttt{\textbackslash g}, \texttt{\textbackslash h}, \texttt{\textbackslash sl}, \texttt{\textbackslash xto}), with extra wrapper packages (bm, bbm, dsfont, xspace, cancel, mathtools). The delivered numbering is the unit's own. Assets: no retained span contains a figure environment, a graphics-inclusion command, a PDF-page inclusion, a table or a TikZ picture; no image file is copied into this package, the package declares no asset dependency and no image file is redistributed with it. Licence: version arXiv:2608.03314v1 of this article is distributed under the Creative Commons Attribution 4.0 International licence, as recorded in the arXiv licence metadata for this exact version and shown on the abstract page https://arxiv.org/abs/2608.03314v1. A full rights notice is delivered at licenses/arxiv-2608.03314-CC-BY-4.0.txt. Characters: every retained excerpt is pure ASCII, so the delivered engine, XeLaTeX with its default Latin Modern text font, reports no missing character.
\begin{theorem}\label{thm}
    Let $G,V$ be one of the above cases. Then we have
    \[\ch_q\C[\N_V]=\sum_{\lambda\in\Lambda_+}\M_q(\lambda)\chi_\lambda.\]
\end{theorem}

\begin{proposition}\label{prop:hesselink}
    Let $V$ be a cofree representation of $G$, and let $P\in\Z[q][t^{-\alpha}]_{\alpha\in S}$ be a function. Let $M_q:\Lambda_+\to\Z[q]$ be a function such that
    \[\sum_{w\in W}(-1)^{\ell(w)}t^{w\rho}\sum_{\lambda\in\Lambda_+}t^{w\lambda}M_q(\lambda)=\sum_{w\in W}(-1)^{\ell(w)}t^{w\rho}w\left(\frac{P}{\prod_{\alpha\in S}(1-qt^{-\alpha})}\right).\]
    Then we have
    \[\ch_q\C[\N_V]=\sum_{\lambda\in\Lambda_+}M_q(\lambda)\chi_\lambda\]
    if and only if
    \begin{equation}
        \frac{\sum_{w\in W}(-1)^{\ell(w)}t^{w\rho}}{\dim_q\C[V]^G}=\sum_{w\in W}(-1)^{\ell(w)}t^{w\rho}wP.\label{eq:hesselink}
    \end{equation}
\end{proposition}

\begin{proposition}\label{prop:2cycle}
    Theorem \ref{thm} holds for the extended quiver representation corresponding to the $2$-cycle.
\end{proposition}

\begin{proof}
    Let $A=\C[V^{\rho^\vee}]^\h$. By Proposition \ref{prop:hesselink}, we wish to show
    \[\sum_{w\in W}(-1)^{\ell(w)}t^{w\rho}\sum_{\lambda\in\Lambda_+}t^{w\lambda}\M_q(\lambda)=\sum_{w\in W}(-1)^{\ell(w)}t^{w\rho}w\left(\frac{1}{\dim_qA\prod_{\alpha\in S_{\leq 0}}(1-qt^{-\alpha})}\right).\]
    Using the definition of $\M_q$, this is equivalent to
    \[\dim_qA\sum_{w\in W}(-1)^{\ell(w)}t^{w\rho}\sum_{\lambda\in\Lambda}\P_q(w\lambda)t^{w\lambda}=\sum_{w\in W}(-1)^{\ell(w)}t^{w\rho}\left(\frac{1}{\prod_{\alpha\in S_{\leq 0}}(1-qt^{-w\alpha})}\right).\]
    Let $m:S_0\to\Z_{\geq 0}$ be a function. We will use the shorthand $t^m$ to mean $\prod_{\alpha\in S_0}t^{m(\alpha)\alpha}$ and $|m|$ to mean $\sum_{\alpha\in S_0}m(\alpha)$. Also, let $WN(S_0)$ be the set of weakly negative functions on $S_0$. We wish to show that
    \begin{equation}\label{eq:2cycle}
        t^\rho\frac{1}{\prod_{\alpha\in S_-}(1-qt^{-\alpha})}\left(\dim_qA\sum_{m\in WN(S_0)}q^{|m|}t^{-m}-\sum_{m\in\Z_{\geq 0}S_0}q^{|m|}t^{-m}\right).
    \end{equation}
    is zero after antisymmetrization. Note that $\sum_{m\in\Z_{\geq 0}S_0\setminus(S_0)_0}q^{|m|}t^{-m}$ is exactly the $q$-character of the space of generators of $\C[V^{\rho^\vee}]$ as an $A$-module, and thus the part of \eqref{eq:2cycle} in parentheses exactly measures the $q$-character of the syzygies of $\C[V^{\rho^\vee}]$ as an $A$-module.

    Let us compute these syzygies. Weights for $\g=\sl_2\oplus\sl_2$ can be described as ordered pairs of integers, and we have $S_0=\{(1,-1),(1,-1),(-1,1),(-1,1)\}$. Let $a_1,a_2,b_1,b_2$ be the generators of $\C[V^{\rho^\vee}]$ such that the $a_i$ have weight $(1,-1)$ and the $b_i$ have weight $(-1,1)$. Then $A$ is generated by
    \[p_1:=a_1b_1,\quad p_2:=a_1b_2,\quad p_3:=a_2b_1,\quad p_4:=a_2b_2\]
    with the relation $p_1p_4=p_2p_3$. The generators of $\C[V^{\rho^\vee}]$ over $A$ are given by $1$, $e_{d,k}:=a_1^{d-k}a_2^k$ for $d>0,0\leq k\leq d$, and $f_{d,k}:=b_1^{d-k}b_2^k$. 

    For $d\in\Z$, let $\C[V^{\rho^\vee}]_d$ be the $A$-submodule of $\C[V^{\rho^\vee}]$ consisting of elements with weight $(d,-d)$. For $d\geq 0$, the generators of this submodule are $e_{d,k}$ for $0\leq k\leq d$. A resolution of $\C[V^{\rho^\vee}]_d$ is given by
    \[\cdots\xto{\partial_2}A\{f^2_k,g^2_k\}_{0\leq k\leq d-1}\xto{\partial_1}A\{f^1_k,g^1_k\}_{0\leq k\leq d-1}\xto{\partial_0}A\{e_{d,k}\}_{0\leq k\leq d},\]
    where
    \begin{align*}
        \partial_0(f_k^1)&=p_3e_{d,k}-p_1e_{d,k+1},&\partial_0(g_k^1)&=p_4e_{d,k}-p_2e_{d,k+1},\\
        \partial_i(f_k^{i+1})&=p_3f_k^i-p_1g_k^i,&\partial_i(g_k^{i+1})&=p_4f_k^i-p_2g_k^i
    \end{align*}
    if $i>0$ is even, and
    
    \[\partial_i(f_k^{i+1})=p_2f_k^i-p_1g_k^i,\quad\partial_i(g_k^{i+1})=p_4f_k^i-p_3g_k^i\]
    if $i$ is odd. Indeed, a relation $F=0$ for $F\in A\{e_{d,k}\}$ can be assumed homogeneous in its multidegree in $b_1,b_2$. Then a straightforward induction argument shows $F$ can be generated by $\partial_0(f_k^1),\partial_0(g_k^1)$, which are the relations with multidegree $(1,0),(0,1)$, respectively. So $\coker\partial_0=\C[V^{\rho^\vee}]_d$. We have that $\ker\partial_0$ is generated by elements in $A\{f_k^1,g_k^1\}$ for fixed $k$. Indeed, this is clear since given $F\in\ker\partial_0$, we can consider which $e_{d,k}$ appear in $\partial_0(F)$. From there, it is easy to see that $\ker\partial_0$ is generated by $\partial_1(f_k^2),\partial_1(g_k^2)$. The same argument then shows exactness at each $A\{f_k^r,g_k^r\}$.
    
    Thus we obtain
    \begin{align*}
        \dim_q\C[V^{\rho^\vee}]_d&=\dim_qA\left((d+1)q^d-2d(q^{d+2}-q^{d+4}+q^{d+6}-\cdots)\right)\\
        &=\dim_qA\left((d+1)q^d-\frac{2dq^{d+2}}{1+q^2}\right)
    \end{align*}
    where the $(d+1)q^d$ term comes from the generators and the rest of the terms from the syzygies. A similar calculation yields the same result when $d<0$.

    Thus we obtain
    \[\dim_qA\sum_{m\in\Z_{\geq 0}S_0\setminus(S_0)_0}q^{|m|}t^{-m}-\sum_{m\in\Z_{\geq 0}S_0}q^{|m|}t^{-m}=\dim_qA\sum_{d\neq 0}t^{(d,-d)}\left(\frac{2|d|q^{|d|+2}}{1+q^2}\right).\]
    We can write
    \[\sum_{d>0}dt^{(d,-d)}q^{d-1}=\frac{t^{(1,-1)}}{(1-qt^{(1,-1)})^2},\quad\sum_{d>0}dt^{(-d,d)}q^{d-1}=\frac{t^{(-1,1)}}{(1-qt^{(-1,1)})^2},\]
    so \eqref{eq:2cycle} can be written as
    \begin{align*}
        &\frac{2q^3\dim_qA}{1+q^2}\left(\frac{t^{(1,1)}}{(1-qt^{(1,1)})^2}\right)\left(\frac{t^{(1,-1)}}{(1-qt^{(1,-1)})^2}+\frac{t^{(-1,1)}}{(1-qt^{(-1,1)})^2}\right)\\
        =&\frac{2q^3\dim_qA}{1+q^2}\left(\frac{t^{(2,0)}}{(1-qt^{(1,1)})^2(1-qt^{(1,-1)})^2}+\frac{t^{(0,2)}}{(1-qt^{(1,1)})^2(1-qt^{(-1,1)})^2}\right).
    \end{align*}
    The action of the two simple reflections in $W$ on weights are via $(a,b)\mapsto(a,-b)$ and $(a,b)\mapsto(-a,b)$. The two terms in the above are invariant under the former and latter of these, respectively. Thus both terms go to zero after antisymmetrization.
\end{proof}

\end{document}
